\%%%%%%%%%%%%%%%%%%%%%%%%%%%%%%%%%%%%%%%%%%%%%%%%%%%%%%%%%%%%%%%%%%%%%%%%%%%%%%%%%%%

\blankpage
\pagebreak
\thispagestyle{empty}
\addcontentsline{toc}{part}{Essaouira}
\mbox{}
\begin{center}

\vspace{3cm}

{\brabo\HUGE\itshape Essaouira}
\end{center}
\pagebreak

%%%%%%%%%%%%%%%%%%%%%%%%%%%%%%%%%%%%%%%%%%%%%%%%%%%%%%%%%%%%%%%%%%%%%%%%%%%%%%%%%%%

\chapter{Skala de la Ville}

A Skala de la Ville é uma das principais estruturas defensivas de Essaouira. Construída junto às muralhas, sua plataforma reúne uma longa fileira de canhões voltados para o Atlântico. O conjunto mostra a combinação entre técnicas de fortificação europeias e arquitetura norte-africana que caracteriza a cidade.

\chapter{Skala do Porto e Bab el-Marsa}
\chapter*{Skala do Porto\\ \& Bab el-Marsa}
\addcontentsline{toc}{chapter}{Skala do Porto \& Bab el-Marsa}

A Skala do Porto protege a entrada do antigo porto e forma, junto à Bab el-Marsa, um dos conjuntos arquitetônicos mais conhecidos de Essaouira. A porta monumental liga a medina à área portuária e concentra alguns dos melhores pontos para observar a atividade do porto. O contraste entre as muralhas, os barcos de pesca e o Atlântico é uma das imagens mais características da cidade.

\chapter{Porto de Essaouira}

O porto foi fundamental para a história econômica da cidade. Desde sua fundação, Essaouira funcionou como um importante porto de comércio internacional, ligando o interior do Marrocos e as rotas caravaneiras ao Atlântico, à Europa e a outras regiões da África. Hoje continua sendo uma área de trabalho, especialmente ligada à pesca. Vale visitar pela manhã, quando a atividade é maior.

\chapter{Medina \& souks}

A medina de Essaouira tem uma organização urbana diferente das medinas medievais de Fès e Marraquexe. Sua construção foi planejada e incorpora ruas mais regulares, eixos comerciais e grandes portas de acesso. Caminhar sem roteiro rígido permite observar lojas, oficinas, pequenos mercados e casas tradicionais.

Os souks são especialmente conhecidos pelo trabalho em madeira de tuia, além de tecidos, cerâmica, objetos de metal, especiarias e artesanato local. A cidade também possui muitas galerias de arte e pequenas lojas voltadas para produção contemporânea.

\chapter{Place Moulay Hassan}

A Place Moulay Hassan fica entre a medina e o porto e funciona como um dos principais espaços de circulação da cidade. É um bom ponto para observar a relação entre o centro urbano, as muralhas e o Atlântico. A praça também concentra cafés e restaurantes.

\chapter{Mellah de Essaouira}

O Mellah é o antigo bairro judaico de Essaouira e ocupa uma posição importante na história comercial da cidade. Desde o século \textsc{xviii}, a comunidade judaica participou ativamente das relações comerciais entre o sultanato e a Europa. Alguns comerciantes judeus receberam o título de \textit{Toujjar es-Sultan}, ou ``comerciantes do sultão'', e desempenharam funções importantes nas redes de comércio e diplomacia.

Há também a Bayt Dakira, um centro dedicado à memória judaica de Essaouira e às relações históricas entre as comunidades judaica e muçulmana. O espaço funciona na antiga sinagoga Simon Attias e reúne elementos relacionados à história religiosa e social da comunidade judaica local. E a sinagoga Chaim Pinto, associada a uma das famílias rabínicas mais importantes da história judaica de Essaouira.

\chapter{Kasbah}

A Kasbah correspondia à área administrativa e militar da cidade. Sua posição junto à medina e às fortificações ajuda a compreender a organização urbana planejada no século \textsc{xviii}. 

\chapter{Musée Sidi Mohammed Ben Abdellah}

O Musée Sidi Mohammed Ben Abdellah funciona em um antigo \textit{riad} e reúne objetos relacionados à história e às tradições culturais de Essaouira e da região. A coleção inclui instrumentos musicais, objetos de artesanato, vestimentas e elementos ligados à vida cotidiana.

\chapter{Litoral e ilhas de Mogador}

A praia de Essaouira se estende ao sul da medina e é um bom lugar para caminhar, especialmente no fim da tarde. A cidade é conhecida pelos ventos constantes, que também explicam sua popularidade entre praticantes de esportes aquáticos.

As ilhas de Mogador ficam diante da cidade e fazem parte do conjunto histórico e natural associado à medina. O arquipélago inclui antigas estruturas defensivas e é também importante para a fauna marinha e de aves. O acesso às ilhas depende das condições do mar e das regras de conservação. Mesmo quando não é possível desembarcar, o arquipélago pode ser observado a partir das muralhas e do \mbox{litoral.}

\chapter{Outros pontos importantes}

Essaouira é um dos principais centros associados à tradição musical Gnawa, manifestação de origem africana incorporada à cultura marroquina. A cidade recebe o Festival Gnaoua e mantém espaços onde é possível ouvir música tradicional ao longo do ano.

Há também o trabalho em madeira de tuia, que é uma das especialidades artesanais mais associadas a Essaouira. A madeira é utilizada em caixas, mesas, objetos decorativos e pequenas peças de marchetaria. Vale observar as oficinas durante a caminhada pela medina, distinguindo os ateliês que produzem as peças localmente das lojas voltadas principalmente ao turismo.